
Consider two agents that both achieve 95\% test pass rate on programming problems. Agent A makes 100 code modifications through trial-and-error; Agent B makes 10 strategic root-cause fixes. Current benchmarks report identical performance---both achieve 95\% accuracy. Yet their debugging processes reveal different capabilities: Agent A iterates blindly through possibilities, while Agent B identifies patterns, clusters errors by shared characteristics, and targets fixes resolving multiple failures simultaneously. Without process metrics, we cannot measure or improve this strategic debugging ability.

The gap becomes apparent when examining multi-test benchmarks. Programming problems with 15+ test cases per problem reveal debugging trajectories invisible to single-attempt evaluation. An agent that clusters three failures sharing ''index out of range'' errors and fixes all three with one modification exhibits behavior distinct from an agent that addresses each failure separately. This strategic capability---iterating efficiently through error feedback---remains unmeasured by existing frameworks.

Existing benchmarks focus on outcomes. HumanEval and MBPP use pass@k to measure generation success rate, evaluating whether agents produce correct code but not how they debug failing code. Competitive programming benchmarks report final pass rates; SWE-bench evaluates GitHub issue resolution without analyzing debugging trajectories when test suites exist. Without process metrics, we cannot distinguish strategic debugging (conceptual understanding enabling root-cause identification) from sequential trial-and-error (iteration without pattern recognition).

This work introduces fix-impact-ratio, a metric measuring tests passed per code modification. Strategic agents achieve ratio > 2.0 by targeting root causes resolving multiple failures simultaneously, while sequential agents show ratio $\approx$ 1.0 (one fix per failure). Controlled experiments validate a two-stage mechanism with large effect sizes: agents cluster errors by type (coefficient 1.909, p=0.001), then prioritize high-impact fixes (42.5\% vs 19.8\%, 2.$15\times$ improvement). However, pattern-based transfer to held-out tests failed (slope ratio 0.82, p=0.504, not significant), revealing strategic debugging operates as a feedback loop effective within revealed test execution, not a learning system enabling pattern transfer to unseen tests without error messages.

Our contributions are: (1) a fix-impact-ratio framework with three execution-based metrics (fix-impact-ratio for efficiency, clustering coefficient for pattern recognition, held-out slope for generalization), (2) experimental validation on controlled mock agents demonstrating the two-stage mechanism with large effect sizes while transfer learning fails clearly, and (3) theoretical clarification that strategic debugging is a feedback loop enabling iterative improvement with error messages, not a learning system supporting predictive generalization without test execution. The framework establishes feasibility of process-based debugging evaluation on controlled data; extension to production systems requires real-world validation.

We organize the paper as follows: Section 2 positions our work relative to code generation benchmarks and execution-feedback approaches. Section 3 describes the fix-impact-ratio framework and metric designs. Section 4 details experimental protocol and hypotheses. Section 5 presents results validating the two-stage mechanism and transfer failure. Section 6 discusses theoretical interpretation, limitations (mock implementation), and future work (real-world validation). Section 7 concludes with implications for agentic evaluation.
